\clearpage
\VerdierPhysicalPageAnchor{0198}%
\begingroup\small\setlength{\parskip}{0pt}\setlength{\emergencystretch}{1.5em}
\setlength{\abovedisplayskip}{4pt plus 2pt minus 2pt}%
\setlength{\belowdisplayskip}{4pt plus 2pt minus 2pt}%

\VerdierSemanticAnchor{ch3d-2-3-1-proof-b-conclusion}
したがって $v$ はエピ射である。これで性質 $(E_3)$ が示され、命題の証明が完了した。

\VerdierSemanticAnchor{ch3d-2-3-2}
\medskip\noindent\textbf{命題 2.3.2.}
\textit{$M$ を $\mathcal A$ の対象の集合で、性質
$(E_1)$、$(E_2)$、$(E_3)$ を満たすものとする。
$Y\in\operatorname{Ob}\operatorname{comp}(\mathcal A)$ を、各整数 $n$ に対して
$H^n(Y)\in M$ であり、かつある整数 $n_0$ が存在して
$n\geq n_0$ なら $H^n(Y)=0$ となる複体とする。このとき、各成分が
$M$ に属する複体
$X\in\operatorname{Ob}\operatorname{comp}^{-}(\mathcal A)$ と擬同型}
\[
\dot u\colon X\longrightarrow Y.
\]
\textit{が存在する。}

\VerdierSemanticAnchor{ch3d-2-3-2-proof}
まず、(1.2.1) により複体
$Y'\in\operatorname{Ob}\operatorname{comp}^{-}(\mathcal A)$ と擬同型
$Y'\longrightarrow Y$ が存在する。したがって、$Y$ が
$\operatorname{comp}^{-}(\mathcal A)$ の対象である場合に帰着できる。
$n$ に関する帰納法により、射の列

\VerdierSemanticAnchor{ch3d-2-3-2-induction}
\[
u_n\colon X_n\longrightarrow X_{n-1}
\]
で次の性質をもつものを構成する：

\VerdierSemanticAnchor{ch3d-2-3-2-induction-properties}
\noindent 1) $n$ が十分小さいとき、$X_n=Y$ かつ
$u_n=\operatorname{id}$ である。

\noindent 2) 複体 $X_n$ の次数 $i\geq -n$ の成分は $M$ に属する。
$X_n$ の次数 $i\geq -n$ の微分の核と像は $M^\wedge$ に属する（2.1.1）。

\noindent 3) 射 $u_n$ は擬同型である。これは射
$p((\dot u_n)^{-n},-n)$（2.1.1.3）と同型である。

\VerdierSemanticAnchor{ch3d-2-3-2-limits}
この構成ができれば、$X=\varprojlim_n X_n$ とおき、標準射
$\dot u$ として次をとればよい：
\[
\dot u\colon X=\varprojlim_n X_n\longrightarrow\varinjlim_n X_n=Y.
\]

\VerdierSemanticAnchor{ch3d-2-3-2-induction-step}
$n$ が十分小さいとき、$X_n=Y$ かつ
$\dot u_n=\operatorname{id}_Y$ とおく。各整数 $n<q$ に対して、
複体 $X_n$ と、上の性質 1)、2)、3) をもつ射 $\dot u_n$ が構成されたと仮定する。
複体 $X_{q-1}$ は次の形をする：

\VerdierSemanticAnchor{ch3d-2-3-2-xqminus1}
\[
X_{q-1}=\cdots\longrightarrow Y^{-q-2}\longrightarrow Y^{-q-1}
\longrightarrow Z^{-q}\longrightarrow L^{-q+1}\longrightarrow L^{-q+2}
\longrightarrow\cdots,
\]
ここで $Y^i$ は複体 $Y$ の成分であり、$L^i$ は $M$ に属する。
このとき、次のエピ射がある：

\VerdierSemanticAnchor{ch3d-2-3-2-quotient}
\[
Z^{-q}\longrightarrow
Z^{-q}/\operatorname{Im}(Y^{-q-1}\longrightarrow Z^{-q})
\longrightarrow0
\]
\endgroup
